Table~\ref{table:pred_examples} lists a few qualitative examples. We discuss next the first example ``How do preschoolers react when caregivers and teachers belittle their autonomous actions'' from the book ``The Whole Child: Development in the Early Years'':

\textbf{Mistral} could correctly guess that the query can be answered by ``2.2.3 Initiative vs. Guilt (Preschool Years)'' using its world knowledge. It doesn't understand that it needs to pick only one option -- due to lack of the task training\\
\textbf{BM25} suffers from the known bias of keyword matching where and picks the section ``8.3.2. Moral Development'' purely because of highest lexical overlap with query tokens, which is the incorrect answer\\
\textbf{DPR} is used out of the box and does a word sensed matching where it predicts ``2.2.2 Autonomy vs. Shame/Doubt (Toddlerhood)'' being most relevant to the query tokens mentioning ``belittling of autonomous actions''. However, without the corpus specific training it falls short of identifying the correct section.\\
\textbf{DSI} is finetuned for the corpus but faulters to find the precise section. It errs just like DPR in this example but might have chosen a more sensible match ``2.2.2 Autonomy vs. Shame/Doubt (Preschool years)'' which takes into account the importance of keyword ``preschool'' for getting a better topic match. However, careful consideration shows this is a hallucinated header not present in the input {\toc}. This correlates with our ablation in figure~\ref{fig:ablation_a} on DSI hallucination.\\
\textbf{\sys}: Our proposed system {\sys} correctly identifies the gold leaf node as ``Initiative vs. Guilt (Preschool Years)''. 

The comparison between DPR, DSI and {\sys} is interesting as they all are built on top odf the same base LLM -- Mistral. The difference between DSI and {\sys} largely points towards the importance of having {\toc} as an input during training, which makes it easier for the LLM to align much better to the corpus. Other anecdotal examples in Table~\ref{table:pred_examples} also enumerate the same partial ordering seen among the baselines and provide more qualitative analysis to support the claim that the design of {\sys} helps it do a much better retrieval.

% We show some anecdotal examples from {\ben} in Table~\ref{table:pred_examples} which further demonstrates the above analysis with real data instances. For example1, we can see how Mistral errs in producing a leaf node showing the lack of task training to generate from {\toc} leaf nodes. Example 2 and example 3 show two different cases where DSI gets it incorrect for different reasons. 

% For example 2, the gold leaf node had only $4$ training examples compared to an average of $18.64$ training examples across all leaf nodes from the same book. As we have already seen in Figure~\ref{fig:ablation_b}, DSI finds it hard to remember the gold leaf node \textit{19.3 Oxidation of Aldehydes and Ketones} possibly due to lack of enough training examples and goes on to produce a hallucinated leaf node.

% For example 3 too,  DSI errs because lack of enough training examples for the gold leaf node \textit{9.5.1 What are non-democracies?}, which is $2$ as compared to an average of $11.79$ training examples for the book. However, DSI instead of hallucinating, predicts a different leaf node \textit{6.2.1 The Experimental Method } which is seemingly relevant and has $21$ training examples for DSI to remember it. 

% For all 3 examples {\sys} produces the right leaf node by learning to utilize the input {\toc} effectively.






% \begin{table}[ht]
% \caption{Predictions for the book \href{https://open.umn.edu/opentextbooks/textbooks/221}{Teaching in a Digital Age}}
% \label{table:pred_examples}
% \begin{center}
% % \addtolength{\tabcolsep}{-0.4em}
% \begin{tabular}{lrl}
% \hline
% \multirow{4}{*}{Ex1}&\textbf{Query}:& How do cultural factors influence discussions in online education? \\
% & \textbf{Gold}:&4.4 Online collaborative learning \\
% & \textbf{Mistral}:&9.10 Deciding (Chapter 9) \\
% & \textbf{DSI}:&12.10 Step eight: communicate, communicate, communicate \\
% & \textbf{{\sys}}:&4.4 Online collaborative learning \\
% \hline
% \multirow{5}{*}{Ex2}&\textbf{Query}:& What measures must be taken to guarantee a gradual advancement \\
% & & in a student's ability over the course of a curriculum?\\
% & \textbf{Gold}:&6.5 Developing skills\\
% & \textbf{Mistral}:&12.6 Step six: set appropriate learning goals.\\
% & \textbf{DSI}:&8.5 Computing\\
% & \textbf{{\sys}}:&6.5 Developing skills\\
% \hline
% \multirow{5}{*}{Ex3}&\textbf{Query}:&What is the suggested time allocation for students studying \\
% & & a course or program?\\
% & \textbf{Gold}:&12.9 Step seven: design course structure and learning activities\\
% & \textbf{Mistral}:&12.6 Step six: set appropriate learning goals\\
% & \textbf{DSI}:&9.3 Ease of Use	\\
% & \textbf{{\sys}}:&12.9 Step seven: design course structure and learning activities\\
% \end{tabular}
% \end{center}
% \end{table}

% We look at some predictions for the book ``Teaching in a Digital Age`` in Table ~\ref{table:pred_examples}. Table of contents for this book is highlighted in Figure~\ref{fig:edu_book3_toc}. In Ex1, Mistral generates incorrect leaf as output and it indicates that it does not understand that it only needs to generate the leaf node, even though it is clearly instructed to do so in the prompt (Section ~\ref{subsec:toc_prompt}). DSI system makes an incorrect prediction as it is not able to leverage the fact that a query about online education could probably be answered by Chapter 4 which is titled ``Methods for teaching with an online focus``. {\sys} makes the correct prediction! In Ex2, DSI outputs a completely unrelated section, as it is unable to leverage the rich information present in Table of Contents -- that ``Chapter 6'' talks about learning and thus ``6.5'' is a better fit for the user query. Ex3 highlights another issue where Mistral model predicts incorrect (hallucinates) leaf node and DSI model predicts a completely unrelated leaf node.